\section{Related Work}
\label{sec:related_work_reframed}

\textbf{Gradient inversion and the role of the observation model.}
Gradient inversion reconstructs a training input by finding a synthetic input whose gradient or update agrees with the observation available to an attacker. DLG introduced direct gradient matching~\cite{zhu2019dlg}; iDLG exploited label information~\cite{zhao2020idlg}; and Inverting Gradients used a cosine objective and total-variation prior to obtain high-quality image reconstructions~\cite{geiping2020inverting}. Subsequent work has demonstrated that batch reconstruction, priors, network design, and the precise local-training mechanism materially change attack performance~\cite{yin2021gradinversion,wei2020framework,du2025sok}. The relevant attack target is therefore not an abstract ``gradient,'' but the actual client-to-server observation: a raw gradient, a multi-step update, a compressed message, or a state dictionary containing non-gradient tensors. This distinction is central here because the implemented FL protocol transmits batch-normalization (BN) running state alongside trainable parameters.

BN has two different roles in the literature. GradInversion uses batch statistics as a reconstruction prior~\cite{yin2021gradinversion}, whereas Hatamizadeh et al. show that updated BN running statistics and the exact global checkpoint must be modeled to invert realistic epoch-wise FL updates~\cite{hatamizadeh2023gia}. Their result is particularly relevant to this paper: an attacker can exploit information in the transmitted state even if a nominal gradient-only optimizer is weak. Our analysis does not claim to discover BN leakage. Instead, it treats all transmitted tensors, including BN state, as part of the server observation and asks whether a proposed transformation changes the leakage available from that observation.

\textbf{Tabular leakage and calibrated baselines.}
Image reconstructions can often be visually inspected, but this is unreliable for mixed categorical--continuous records. TabLeak addresses this problem with a structured tabular attack, ensemble pooling, a feature-level accuracy metric, and a random marginal baseline~\cite{vero2023tableak}. On Adult, TabLeak reports that a result below the marginal baseline does not establish extraction of private information. This is a direct precedent for our use of input-space, data-free baselines, and it rules out any claim that such baseline calibration is new. We use TabLeak instead as a native positive-control instrument for tabular FL, paired with an image-domain Geiping control, so that a failure of a generic attacker is not misread as evidence for a defense.

\textbf{Defenses, adaptive attacks, and audit methodology.}
Gradient pruning, representation perturbation, input encoding, clipping, sparsification, and noise have all been proposed as ways to weaken inversion. Huang et al. systematically evaluated gradient pruning, MixUp, and Intra-InstaHide under a strong attacker that knows labels and private-batch BN statistics~\cite{huang2021evaluating}. They emphasized accuracy--leakage tradeoffs but left direct DP-SGD comparison to future work. Balunovi\'{c} et al. formalized gradient leakage through a Bayesian observation model and showed why the defense distribution must enter the adversary's inference problem~\cite{balunovic2022bayesian}. Yue et al. then showed that gradient obfuscation defenses, including representation-based and stochastic mechanisms, can be circumvented by an attack adapted to the post-processing~\cite{yue2023obfuscation}. Li et al. similarly audit additive noise, clipping, sparsification, and Soteria with a generative leakage attack~\cite{li2022auditing}. These papers make clear that a non-adaptive or under-tuned attack cannot be treated as a privacy certificate.

The recent SoK extends this perspective to practical FL settings, evaluating quantization, sparsification, clipping, and perturbation jointly with utility and emphasizing the fragility of GIA conclusions under training and architectural choices~\cite{du2025sok}. Thus, our contribution is not the first systematization of gradient-leakage defenses, the first adaptive attack, or the first privacy--utility comparison. It is an operational audit protocol for a selected heterogeneous defense suite. Before interpreting a defense result, we require: (i) a successful undefended positive control; (ii) an attacker given exactly the information held by the server; (iii) input-space scoring that beats a data-free baseline; and (iv) a DP comparator whose clipping scope is specified and whose operating points are matched both by update distortion and by downstream utility. Existing studies contain important subsets of these checks, but generally not their conjunction across both the image and tabular native attack settings considered here.

\textbf{Keyed transforms, aggregation, and DP.}
Secret randomness is not itself a DP guarantee. InstaHide demonstrated that an encoding can preserve task accuracy while concealing raw-looking inputs~\cite{huang2020instahide}, but subsequent cryptanalysis showed that instance encoding without suitable formal protection can fail to provide meaningful privacy~\cite{carlini2021instance}. Random projections and Hadamard rotations are also established communication primitives for distributed mean estimation~\cite{suresh2017dme}; adding calibrated discrete Gaussian or Skellam noise can yield formal DP in appropriate federated protocols~\cite{kairouz2021ddg,agarwal2021skellam}. In contrast, the two seeded transforms studied here make no DP claim. Their evaluation must use the realistic threat model: if the server possesses the seed needed for aggregation, inversion, or recovery of a transmitted statistic, then that seed belongs in the attacker's knowledge. Count-sketch aggregation is treated similarly: its privacy implication depends on whether the server knows the hashes and on which individual or aggregate message is visible.

Accordingly, this paper positions seeded transforms as empirical update transformations, not encryption and not DP. We report their leakage only after passing the same native positive controls used for the published defenses, and we compare them with a specified Gaussian baseline rather than treating a large, uncalibrated noise vector as a meaningful DP reference. This framing aligns the claims with the observation channel actually available to the server and keeps the conclusion limited to the evaluated protocol, models, datasets, and attacker knowledge.
